\subsection{模的同型分量}

定义 4. —— 设 A 是一个环，M 是一个 A-模，S 是一个单 A-模。M 的 S 型同型分量记作 \(M_{S}\)，它是 M 中同构于 S 的子模之和。

显然，\(M_{S}\) 是 M 的 S 型同型的最大子模。由于 \(M_{S}\) 的每个子模都是 S 型同型模（VIII, 第 61 页，命题 2），故对 M 的每个子模 N 都有 \(N_{S} = M_{S} \cap N\)。

若 \(S'\) 是同构于 S 的单 A-模，则显然有 \(M_{S} = M_{S'}\)，因而 \(M_{S}\) 只依赖于 S 的类（VIII, 第 51 页）。

设 M 是一个 A-模。M 存在一个最大的半单子模，称为 M 的底座；它是 M 的诸单子模之和，也是 M 的诸同型分量之和。特别地，M 是半单模当且仅当它等于自身的底座。

命题 4. —— 设 A 是一个环。用 S 表示单 A-模的类组成的集合。设 M 是一个半单 A-模。

\begin{enumerate}
\def\labelenumi{\alph{enumi})}
\item
  模 M 是其同型分量族 \((\mathrm{M}_{\lambda})_{\lambda\in\mathcal{S}}\) 的直和。
\item
  设 M 是单子模族 \((\mathrm{N}_{i})_{i\in\mathrm{I}}\) 的直和。对每个 \(\lambda\in\mathscr{S}\)，用 \(\mathrm{I}(\lambda)\) 表示这样的指标 \(i\in I\) 组成的集合：\(N_{i}\) 属于类 \(\lambda\)。我们有 \(M_{\lambda}=\bigoplus_{i\in\mathrm{I}(\lambda)}N_{i}\)。
\item
  若 M 是有限生成的，则诸 \(\lambda\in S\) 中使 \(M_{\lambda}\neq0\) 者组成的集合是有限的。
\item
  设 \(\mathrm{N}\) 是 \(\mathrm{M}\) 的任意一个子模。对每个 \(\lambda \in \mathcal{S}\)，我们有 \(\mathrm{N}_{\lambda} = \mathrm{N} \cap \mathrm{M}_{\lambda}\) 与 \((\mathrm{M} / \mathrm{N})_{\lambda} = (\mathrm{M}_{\lambda} + \mathrm{N}) / \mathrm{N}\)。
\end{enumerate}

由于 M 是半单的，它是族 \((\mathrm{M}_{\lambda})_{\lambda\in\mathcal{S}}\) 之和；我们来证明这个和是直和。设 \(\lambda\in S\)。把族 \((\mathrm{M}_{\mu})_{\mu\in\mathcal{S}-\{\lambda\}}\) 之和记为 \(M_{\lambda}^{\prime}\)。模 \(M_{\lambda}^{\prime}\) 是一族不同构于 \(\lambda\) 的单模的直和（VIII, 第 56 页，定理 1）。由 VIII, 第 56 页定理 1 的推论 3，\(M_{\lambda}^{\prime}\) 不含任何类为 \(\lambda\) 的单子模。因此我们有 \(M_{\lambda}\cap M_{\lambda}^{\prime}=0\)。断言 a) 于是得证。按构造，我们有 \(M_{\lambda}\supset\bigoplus_{i\in\mathrm{I}(\lambda)}N_{i}\)，故断言 b) 由 II, §1, No.~8, 第 208 页的注 1 得出。

断言 c) 由 a) 及 II, §1, No.~12, 第 221 页的命题 23 得出。

设 N 是 M 的一个子模，\(\lambda\in S\)。N 的同型分量 \(N_{\lambda}\)（\(\lambda\) 型）包含于 \(M_{\lambda}\)，且 \(M_{\lambda}\cap N\subset N_{\lambda}\)。因此，交 \(N\cap M_{\lambda}\) 就是 N 的 \(\lambda\) 型同型分量。

对每个 \(\lambda \in \mathcal{S}\)，模 \(\mathrm{M}_{\lambda} + \mathrm{N} / \mathrm{N}\) 同构于 \(\mathrm{M}_{\lambda} / (\mathrm{M}_{\lambda} \cap \mathrm{N})\)。因而它是 \(\lambda\) 型同型模，且包含于 \((\mathrm{M} / \mathrm{N})_{\lambda}\)。最后的断言于是由 a) 及 II, §1, No.~8, 第 208 页的注 1 得出。

推论. —— 设 A 是一个环，S 是单 A-模的类组成的集合。设 M 是一个半单 A-模，N 是 M 的一个子模。则我们有 \(N = \bigoplus_{\lambda \in \mathscr{S}} N \cap M_{\lambda}\) 与 \(M/N = \bigoplus_{\lambda \in \mathscr{S}} (M_{\lambda} + N)/N\)。

由于 N 与 M/N 都是半单的（VIII, 第 56 页，推论 3），本推论由命题 4, d) 得出。

半单 A-模 M 的支集是使 M 的 \(\lambda\) 型同型分量非零的单 A-模的类 \(\lambda\) 组成的集合。有限生成的半单 A-模的支集是有限的。

命题 5. —— 设 A 是一个环，S 是单 A-模的类组成的集合。设 M 与 N 是 A-模。

\begin{enumerate}
\def\labelenumi{\alph{enumi})}
\item
  设 \(f: \mathrm{M} \to \mathrm{N}\) 是一个同态。对每个 \(\lambda \in \mathcal{S}\)，\(f\) 诱导出一个同态 \(f_{\lambda}\)（从 \(\mathrm{M}_{\lambda}\) 到 \(\mathrm{N}_{\lambda}\)）；若 \(\mathrm{M}\) 是半单的且 \(f\) 是满射，则每个同态 \(f_{\lambda}\) 都是满射。
\item
  设 M 是半单的。映射 \(f \mapsto (f_{\lambda})_{\lambda \in \mathcal{S}}\) 是从 \(\mathrm{Hom}_{\mathrm{A}}(\mathrm{M},\mathrm{N})\) 到 \(\prod_{\lambda \in \mathcal{S}}\mathrm{Hom}_{\mathrm{A}}(\mathrm{M}_{\lambda},\mathrm{N}_{\lambda})\) 的群同构。当 M 等于 N 时，该映射是从 \(\mathrm{End}_{\mathrm{A}}(\mathrm{M})\) 到 \(\prod_{\lambda \in \mathcal{S}}\mathrm{End}_{\mathrm{A}}(\mathrm{M}_{\lambda})\) 的环同构。
\end{enumerate}

对每个 \(\lambda \in \mathcal{S}\)，子模 \(f(\mathrm{M}_{\lambda})\)（\(\mathbf{N}\) 的）同构于某个 \(\lambda\) 型同型模的商；因而它是 \(\lambda\) 型同型模，从而包含于 \(\mathrm{N}_{\lambda}\)。

设 \(\mathrm{M}\) 是半单的且 \(f\) 是满射。于是 \(f\) 诱导出从 \(\mathrm{M} / \mathrm{Ker}(f)\) 到 \(\mathrm{N}\) 的同构，它把 \((\mathrm{M}_{\lambda} + \mathrm{Ker}(f)) / \mathrm{Ker}(f)\) 映到 \(f(\mathrm{M}_{\lambda})\)。

由 VIII, 第 65 页的命题 4，我们有 \(N_{\lambda} = f(M_{\lambda})\)，这就完成了 a) 的证明。

b) 中所考虑的映射显然是群同态，且当 M 等于 N 时是环同态。设 \((f_{\lambda})_{\lambda \in \mathcal{S}}\) 是 \(\prod_{\lambda \in \mathcal{S}} \mathrm{Hom}(M_{\lambda}, N_{\lambda})\) 的一个元素。它在 b) 中映射下的逆像的唯一元素是由下式定义的同态 \(f : M \to N\)：

\[
f \Bigl (\sum_ {\lambda \in \mathcal {S}} x _ {\lambda} \Bigr) = \sum_ {\lambda \in \mathcal {S}} f _ {\lambda} (x _ {\lambda})
\]

其中 \((x_{\lambda})_{\lambda \in \mathcal{S}}\in \bigoplus_{\lambda}\mathrm{M}_{\lambda}\) 任意。

注. —— 设 A 与 B 是环。设 M 是一个 (A, B)-双模。由命题 5 可知，A-模 M 的诸同型分量都是 M 的子双模。特别地，当 M 是一个 A-模且 B 是 \(\operatorname{End}_{\mathrm{A}}(\mathrm{M})\) 的反环时就属于这种情形。

例. —— 我们来考虑环 A 交换的情形。把极大理想 m 映到 \(\operatorname{cl}(A/m)\) 的映射是从 A 的极大理想组成的集合到单 A-模的类组成的集合 S 的一个双射（VIII, 第 51 页）。逆双射把 \(\lambda\) 映到它的零化理想 \(m_{\lambda}\)。

设 M 是一个 A-模。对每个 \(\lambda\in S\)，M 的同型分量 \(M_{\lambda}\)（\(\lambda\) 型）由被 \(m_{\lambda}\) 零化的元素组成，并且我们可以把 \(M_{\lambda}\) 看作体 \(A/m_{\lambda}\) 上的一个向量空间。若 M 是半单的而 N 是另一个 A-模，则由命题 5 可以得到从 \(\operatorname{Hom}_{\mathrm{A}}(\mathrm{M},\mathrm{N})\) 到 \(\prod_{\lambda\in\mathcal{S}}\operatorname{Hom}_{\mathrm{A}/m_{\lambda}}(\mathrm{M}_{\lambda},\mathrm{N}_{\lambda})\) 的一个群同构。

推论 1. —— 设 M 是一个半单 A-模，N 是 M 的一个子模。下列性质等价：

\begin{enumerate}
\def\labelenumi{(\roman{enumi})}
\item
  在 M 中存在唯一的与 N 互补的子模。
\item
  我们有 \(\mathrm{Hom}_{\mathrm{A}}(\mathrm{M} / \mathrm{N},\mathrm{N}) = 0\)
\item
  存在一个子集 \(\Lambda\)（含于 \(\mathcal{S}\)），使得 \(N = \bigoplus_{\lambda \in \Lambda} M_{\lambda}\)。
\end{enumerate}

选取 M 中与 N 互补的一个子模 \(N^{\prime}\)（VIII, 第 56 页，推论 2）。若我们把 M 与 \(N^{\prime} \times N\) 等同起来，则 M 中与 N 互补的子模就是从 \(N^{\prime}\) 到 N 的 A-线性映射的图象。由于 \(N^{\prime}\) 同构于 M/N，我们便证明了性质 (i) 与 (ii) 等价。由命题 5, b)，群 \(\operatorname{Hom}_{\mathrm{A}}(\mathrm{N}^{\prime}, \mathrm{N})\) 同构于群 \(\prod_{\lambda \in \mathcal{S}} \operatorname{Hom}_{\mathrm{A}}(\mathrm{N}_{\lambda}^{\prime}, \mathrm{N}_{\lambda})\)。它为零当且仅当对每个 \(\lambda \in S\) 有 \(N_{\lambda} = 0\) 或 \(N_{\lambda}^{\prime} = 0\)（VIII, 第 61 页，注），亦即 \(N_{\lambda} = 0\) 或 \(N_{\lambda} = M_{\lambda}\)。这就证明了性质 (ii) 与 (iii) 等价。

推论 2. —— 设 M 是一个 A-模。下列两个条件等价：

\begin{enumerate}
\def\labelenumi{(\roman{enumi})}
\item
  M 的每个子模都有唯一的补子模。
\item
  模 M 是一族两两不同构的单模 \((\mathrm{S}_{i})_{i\in\mathrm{I}}\) 的直和。
\end{enumerate}

设 M 满足这些条件。于是，对 M 的每个子模 N，存在 I 的唯一子集 J 使我们有 \(N = \bigoplus_{j \in J} S_j\)，且 M 的每个单子模都等于某个 \(S_i\)。

条件 (i) 与 (ii) 都蕴含 M 是半单的。

设条件 (i) 成立。设 \(\lambda \in \mathcal{S}\)。由推论 1 中 (i) 与 (iii) 的等价性，\(\mathrm{M}_{\lambda}\) 的每个子模或为零或等于 \(\mathrm{M}_{\lambda}\)。因此 \(\mathrm{M}_{\lambda}\) 或为零或为单模，而 (ii) 由 \(\mathrm{M}\) 是族 \((\mathrm{M}_{\lambda})_{\lambda \in \mathcal{S}}\) 的直和这一事实得出。

反之，若条件 (ii) 成立，则 \(M_{\lambda}\) 对每个 \(\lambda \in S\) 均为零或单模。若 N 是 M 的一个子模，则 \(N \cap M_{\lambda}\) 等于 0 或 \(M_{\lambda}\)（对每个 \(\lambda \in S\)）。由于我们有 \(N = \bigoplus_{\lambda \in S} (N \cap M_{\lambda})\)（VIII, 第 66 页，推论），子模 N 具有推论 1 的性质 (iii)，从而在 M 中有唯一的补子模。这就证明了 (ii) 蕴含 (i)，同时也证明了本推论的各条断言。

设 M 是一个 A-模，S 是一个单 A-模。用 D 表示体 \(\operatorname{End}_{\mathrm{A}}(\mathrm{S})\) 的反环，并把 S 视为 (A,D)-双模。于是 \(\operatorname{Hom}_{\mathrm{A}}(\mathrm{S},\mathrm{M})\) 是 D 上的左向量空间，而 \(\operatorname{Hom}_{\mathrm{A}}(\mathrm{M},\mathrm{S})\) 是 D 上的右向量空间。左 D-向量空间 \(\operatorname{Hom}_{\mathrm{A}}(\mathrm{S},\mathrm{M})\) 的对偶是 D 上的右向量空间（II, §2, No.~3, 第 232 页，定义 2）。对每个 \(u \in \operatorname{Hom}_{\mathrm{A}}(\mathrm{M},\mathrm{S})\)，映射 \(h(u): v \mapsto u \circ v\)（从 \(\operatorname{Hom}_{\mathrm{A}}(\mathrm{S},\mathrm{M})\) 到 \(\operatorname{Hom}_{\mathrm{A}}(\mathrm{S},\mathrm{S}) = \mathrm{D}\)）是左 D-向量空间 \(\operatorname{Hom}_{\mathrm{A}}(\mathrm{S},\mathrm{M})\) 上的线性型。

命题 6. —— 沿用上面的记号，并设 A-模 M 是半单的。映射 \(u \mapsto h(u)\)（从右 D-向量空间 \(\mathrm{Hom}_{\mathrm{A}}(\mathrm{M},\mathrm{S})\) 到左 D-向量空间 \(\mathrm{Hom}_{\mathrm{A}}(\mathrm{S},\mathrm{M})\) 的对偶）是 D-线性的双射。

设 \(u \in \mathrm{Hom}_{\mathrm{A}}(\mathrm{M}, \mathrm{S})\)，\(v \in \mathrm{Hom}_{\mathrm{A}}(\mathrm{S}, \mathrm{M})\)，\(d \in \mathrm{D}\)。我们有

\[
h (u d) (v) = h (d \circ u) (v) = d \circ u \circ v = d \circ (h (u) (v)) = h (u) (v) d.
\]

这就证明映射 h 是 D-线性的。它就是由 \(u \mapsto \operatorname{Hom}(1_{\mathrm{S}}, u)\) 给出的从 \(\operatorname{Hom}_{\mathrm{A}}(\mathrm{M}, \mathrm{S})\) 到 \(\operatorname{Hom}_{\mathrm{D}}(\operatorname{Hom}_{\mathrm{A}}(\mathrm{S}, \mathrm{M}), \operatorname{Hom}_{\mathrm{A}}(\mathrm{S}, \mathrm{S}))\) 的映射。要证明它是双射，由 VIII, 第 66 页的命题 5, b)，只需处理 M 是 S 型同型模的情形；这时可以应用 VIII, 第 65 页的推论。

